\begin{resumo}[\protect\bfseries Resumo]
  Este trabalho busca detalhar o processo de vivência, bem como os projetos realizados em cada período da AGES (Agência Experimental de Engenharia de Software), objetivando documentar o percurso de formação acadêmica trilhado pela aluna Mariah Backes Freire. A AGES é uma disciplina curricular obrigatória para o curso de Bacharelado em Engenharia de Software que visa promover a atuação prática dos alunos em projetos reais de sua escolha. Cada nível ordenado da AGES aqui documentado (AGES I, II, III e IV) corresponde a um semestre de participação ativa na disciplina, cada um com suas respectivas responsabilidades: programação e testes; programação e projeto de banco de dados; verificação e validação de software, projeto e arquitetura de software; gerenciamento de projetos, aprofundamento de outras competências desenvolvidas no curso e portfólio de conclusão de curso.
  \bigbreak\
  \\\textbf{PALAVRAS CHAVES:}
  AGES, Engenharia de Software, Tecnologia da Informação, desenvolvimento de software, memorial descritivo.
\end{resumo}